\documentclass[11pt,a4paper]{article}
\usepackage[utf8]{inputenc}
\usepackage[T1]{fontenc}
\usepackage[polish]{babel}
\usepackage{amsmath}
\usepackage{amssymb}
\usepackage[margin=1.8cm]{geometry}
\usepackage{lmodern}
\usepackage{setspace}
\pagestyle{empty}
\setlength{\parindent}{0pt}
\setlength{\parskip}{0.8em}
\linespread{1.15}

\begin{document}

{\LARGE\textbf{Pytanie 1}}

Dla zadanej trajektorii:

\[
\vec r(t) = \bigl(2t^2 - t, t^3 - 3t\bigr)
\]

wektor przyspieszenia w chwili $t=2$ wynosi:

\vspace{0.5em}

{\large
\textbf{A.} $(6, 12)$

\textbf{B.} $(4, 12)$

\textbf{C.} $(6, 6)$

\textbf{D.} $(8, 12)$
}

\newpage

{\LARGE\textbf{Pytanie 2}}

Dla zadanej prędkości i warunku początkowego:

\[
\vec v(t) = \bigl(3t^2, 2t\bigr), \quad \vec r(0)=(0,0)
\]

położenie w chwili $t=2$ wynosi:

\vspace{0.5em}

{\large
\textbf{A.} $(8, 4)$

\textbf{B.} $(6, 4)$

\textbf{C.} $(8, 2)$

\textbf{D.} $(4, 2)$
}

\newpage

{\LARGE\textbf{Pytanie 3}}

Dla zadanego przyspieszenia i warunku początkowego:

\[
\vec a(t) = \bigl(6t, -4\bigr), \quad \vec v(0)=(0,0)
\]

prędkość w chwili $t=2$ wynosi:

\vspace{0.5em}

{\large
\textbf{A.} $(12, -8)$

\textbf{B.} $(6, -8)$

\textbf{C.} $(12, -4)$

\textbf{D.} $(6, -4)$
}

\newpage

{\LARGE\textbf{Pytanie 4}}

Dla zadanej trajektorii oraz masy ciała:

\[
\vec r(t) = \bigl(t^3 - 2t, 2t^2 + t\bigr), \quad m = 2
\]

siła działająca na ciało w chwili $t=4$ wynosi:

\vspace{0.5em}

{\large
\textbf{A.} $(48, 8)$

\textbf{B.} $(24, 4)$

\textbf{C.} $(48, 4)$

\textbf{D.} $(24, 8)$
}

\newpage

{\LARGE\textbf{Pytanie 5}}

Dla zadanej trajektorii oraz siły zależnej od położenia:

\[
x(t) = t^2, \quad t \in [0,1]
\]

\[
F(x) = 2x
\]

praca wykonana przez tę siłę w pierwszej sekundzie ruchu wynosi:

\vspace{0.5em}

{\large
\textbf{A.} $\frac{1}{2}$

\textbf{B.} $\frac{2}{3}$

\textbf{C.} $1$

\textbf{D.} $\frac{3}{2}$
}

\newpage

{\LARGE\textbf{Pytanie 6}}

Dla zadanej trajektorii:

\[
\vec r(t) = \bigl(\cos(t^2), \sin(t^2)\bigr), \qquad t \ge 0
\]

które z poniższych stwierdzeń jest prawdziwe?

\vspace{0.5em}

{\large
\textbf{A.} Ruch odbywa się po okręgu, a szybkość jest stała

\textbf{B.} Ruch odbywa się po okręgu, a szybkość rośnie liniowo w czasie

\textbf{C.} Ruch nie odbywa się po okręgu, a szybkość jest stała

\textbf{D.} Ruch nie odbywa się po okręgu, a szybkość rośnie liniowo w czasie
}

\newpage

{\LARGE\textbf{Pytanie 7}}

Piłeczka tenisowa zostaje puszczona z wysokości 4. piętra. Przyjmujemy, że jedno piętro ma wysokość $3\ \text{m}$. Po każdym odbiciu piłeczka traci $\frac{1}{4}$ swojej energii mechanicznej. Na jaką maksymalną wysokość wzniesie się piłeczka po trzecim odbiciu?

\vspace{0.5em}

{\large
\textbf{A.} $12\cdot\left(\frac{3}{4}\right)^2$

\textbf{B.} $12\cdot\left(\frac{3}{4}\right)^3$

\textbf{C.} $12\cdot\left(\frac{1}{4}\right)^2$

\textbf{D.} $12\cdot\left(\frac{1}{4}\right)^3$
}

\newpage

{\LARGE\textbf{Pytanie 8}}

Dla oscylatora harmonicznego o początkowym maksymalnym wychyleniu $A$, który przy każdym przejściu przez położenie równowagi traci $\frac{1}{4}$ swojej energii mechanicznej, maksymalne wychylenie osiągnięte \textbf{bezpośrednio po drugim przejściu przez położenie równowagi} wynosi:

\vspace{0.5em}

{\large
\textbf{A.} $A\left(\frac{3}{4}\right)^2$

\textbf{B.} $A\sqrt{\frac{3}{4}}$

\textbf{C.} $\frac{3}{4}A$

\textbf{D.} $\frac{1}{2}A$
}

\newpage

{\LARGE\textbf{Pytanie 9}}

Dla zadanej fali:

\[
y(x,t)=2\sin(\pi x-4\pi t)
\]

które z poniższych stwierdzeń jest prawdziwe?

\vspace{0.5em}

{\large
\textbf{A.} Fala rozchodzi się z prędkością $4$, a punkty $x=0$ i $x=2$ drgają w tej samej fazie

\textbf{B.} Fala rozchodzi się z prędkością $2$, a punkty $x=0$ i $x=2$ drgają w tej samej fazie

\textbf{C.} Fala rozchodzi się z prędkością $2$, a punkty $x=0$ i $x=2$ drgają w przeciwnej fazie

\textbf{D.} Fala rozchodzi się z prędkością $4$, a punkty $x=0$ i $x=2$ drgają w przeciwnej fazie
}

\newpage

{\LARGE\textbf{Pytanie 10}}

Dane są dwa drgania:

\[
y_1(t)=A\sin(2\pi f_1 t), \qquad y_2(t)=A\sin(2\pi f_2 t),
\]

gdzie $f_1=10\ \text{Hz}$ oraz $f_2=12\ \text{Hz}$.

Częstotliwość dudnień wynosi:

\vspace{0.5em}

{\large
\textbf{A.} $1\ \text{Hz}$

\textbf{B.} $2\ \text{Hz}$

\textbf{C.} $11\ \text{Hz}$

\textbf{D.} $22\ \text{Hz}$
}

\end{document}
